% =====================================================================
%  Espalhamento Profundamente Inelástico — apostila do HADROS Sandbox
%  Compilar com:  latexmk -pdf dis.tex     (pdflatex + biber)
% =====================================================================
\documentclass[11pt,a4paper,oneside]{report}

% ---------- tipografia -------------------------------------------------
% Palatino (Hermann Zapf, 1949) via newpx, com fonte matemática casada.
\usepackage[T1]{fontenc}
\usepackage[utf8]{inputenc}
\usepackage[brazil]{babel}
\usepackage{newpxtext}
\usepackage[varg,bigdelims]{newpxmath}
\usepackage[scaled=0.86]{helvet}
\usepackage[scaled=0.83]{beramono}
\usepackage{microtype}
\linespread{1.06}

% ---------- geometria --------------------------------------------------
\usepackage[a4paper,left=30mm,right=48mm,top=28mm,bottom=30mm,
            marginparwidth=34mm,marginparsep=6mm,footskip=14mm]{geometry}

% ---------- matemática -------------------------------------------------
\usepackage{amsmath}
\let\Bbbk\relax          % newpxmath já define \Bbbk; evita o conflito com amssymb
\usepackage{amssymb,mathtools}
\usepackage{slashed}
\usepackage{bm}
\allowdisplaybreaks[1]

% ---------- gráficos e cor --------------------------------------------
\usepackage{graphicx}
\usepackage[dvipsnames,table]{xcolor}
\usepackage{tikz}
\usetikzlibrary{arrows.meta,decorations.pathmorphing,decorations.markings,
                decorations.pathreplacing,positioning,calc,shapes.geometric,
                fit,backgrounds}
\usepackage{booktabs}
\usepackage[font=small,labelfont=bf]{caption}
\usepackage{subcaption}

% ---------- estrutura --------------------------------------------------
\usepackage{titlesec}
\usepackage{enumitem}
\usepackage[most]{tcolorbox}
\usepackage{fancyhdr}
\usepackage{marginnote}
\usepackage{csquotes}
\usepackage[backend=biber,style=numeric-comp,sorting=none,giveninits=true,
            maxbibnames=6,minbibnames=3,doi=false,isbn=false,url=false,
            eprint=true]{biblatex}
\addbibresource{dis.bib}
\usepackage[hidelinks]{hyperref}
\usepackage[nameinlink,brazilian]{cleveref}

% ---------- paleta -----------------------------------------------------
\definecolor{tinta}    {HTML}{16202B}
\definecolor{acento}   {HTML}{1D4E89}
\definecolor{acento2}  {HTML}{B14A20}
\definecolor{verdemar} {HTML}{157F63}
\definecolor{papel}    {HTML}{FBFAF7}
\definecolor{filete}   {HTML}{C9CED6}
\definecolor{suave}    {HTML}{6B7480}

\pagecolor{papel}
\color{tinta}

% ---------- títulos ----------------------------------------------------
\titleformat{\chapter}[display]
  {\normalfont\sffamily}
  {\raggedright\large\bfseries\color{acento}\MakeUppercase{Capítulo \thechapter}}
  {0.6ex}
  {\color{tinta}\Huge\bfseries\rmfamily\raggedright}
  [\vspace{1.2ex}{\color{filete}\titlerule[1.1pt]}]
\titlespacing*{\chapter}{0pt}{-8mm}{22pt}

\titleformat{\section}
  {\normalfont\sffamily\large\bfseries\color{acento}}{\thesection}{0.8em}{}
\titleformat{\subsection}
  {\normalfont\sffamily\normalsize\bfseries\color{tinta}}{\thesubsection}{0.7em}{}
\titlespacing*{\section}{0pt}{2.6ex plus 1ex minus .2ex}{1.2ex plus .2ex}

% ---------- cabeçalhos -------------------------------------------------
\pagestyle{fancy}
\fancyhf{}
\renewcommand{\headrulewidth}{0pt}
\fancyhead[L]{\footnotesize\sffamily\color{suave}\nouppercase{\leftmark}}
\fancyfoot[C]{\footnotesize\sffamily\color{suave}\thepage}
\fancypagestyle{plain}{\fancyhf{}\renewcommand{\headrulewidth}{0pt}
  \fancyfoot[C]{\footnotesize\sffamily\color{suave}\thepage}}
\renewcommand{\chaptermark}[1]{\markboth{\thechapter\ \ #1}{}}

% ---------- caixas -----------------------------------------------------
\tcbset{
  caixabase/.style={
    enhanced, breakable, boxrule=0pt, frame hidden, sharp corners,
    left=3.4mm, right=3.4mm, top=2.6mm, bottom=2.6mm,
    borderline west={1.2pt}{0pt}{#1},
    colback=#1!5!papel,
    fonttitle=\sffamily\bfseries\small,
    coltitle=#1, attach title to upper={\par\smallskip},
    % quando a caixa quebra de página, a continuação ganha um rótulo próprio
    % em vez de aparecer como um trecho órfão sem título
    title after break={\normalfont\sffamily\itshape\small (continuação)}
  }
}

\newtcolorbox{definicao}[1][]{caixabase=acento, title={#1}}
\newtcolorbox{deducao}[1][]{caixabase=verdemar, title={#1}}
\newtcolorbox{cuidado}[1][]{caixabase=acento2, title={#1}}
\newtcolorbox{verifique}[1][]{caixabase=verdemar, title={#1}}

\newtcolorbox{historia}[1][]{
  enhanced, breakable, boxrule=0pt, frame hidden, sharp corners,
  colback=papel, left=6mm, right=4mm, top=2mm, bottom=2mm,
  borderline west={0.6pt}{0pt}{filete},
  fontupper=\small\itshape, title={#1},
  fonttitle=\sffamily\bfseries\small\upshape, coltitle=suave,
  attach title to upper={\par\smallskip}
}

% ---------- atalhos ----------------------------------------------------
\newcommand{\md}{\mathrm{d}}
\newcommand{\gev}{\,\mathrm{GeV}}
\newcommand{\cm}{\,\mathrm{cm}}
\newcommand{\Qsq}{Q^{2}}
\newcommand{\xB}{x}
\newcommand{\nota}[1]{\marginnote{\footnotesize\sffamily\color{suave}#1}}
% barra de Feynman: o pacote slashed posiciona a barra corretamente
% em cada tamanho, o que \not sozinho não faz.
\newcommand{\sla}[1]{\ensuremath{\slashed{#1}}}

\newcommand{\fig}[1]{\includegraphics{figuras/#1}}

% =====================================================================
\begin{document}
\pagenumbering{roman}   % a classe report não tem \frontmatter/\mainmatter

% ---------------------------------------------------------------- capa
\begin{titlepage}
\thispagestyle{empty}
\centering
\vspace*{22mm}

{\color{filete}\rule{\textwidth}{1.1pt}}

\vspace{9mm}
{\sffamily\large\color{acento}\MakeUppercase{HADROS Sandbox\quad\textbullet\quad Apostila 01}}

\vspace{12mm}
{\Huge\bfseries Espalhamento\\[2mm] Profundamente Inelástico\par}

\vspace{7mm}
{\Large\itshape Como se enxerga dentro de um próton,\\[1mm]
 e por que isso decide o destino de um neutrino\\[1mm] de energia ultra-alta\par}

\vspace{9mm}
{\color{filete}\rule{\textwidth}{1.1pt}}

\vspace{16mm}
\begin{minipage}{0.80\textwidth}
\centering\small\itshape\color{suave}
Notas de aula para estudantes de pós-graduação que vão trabalhar com o HADROS3.
Deduções completas, sem pular passos; dados reais do HERA; e todas as figuras
reprodutíveis a partir do código que acompanha o texto.
\end{minipage}

\vfill

\begin{minipage}{0.80\textwidth}
\centering\footnotesize\sffamily\color{suave}
Projeto HADROS3\\[1mm]
\texttt{sandbox/apostila/dis.tex}\\[3mm]
\today
\end{minipage}

\vspace{12mm}
\end{titlepage}

% ------------------------------------------------------------- prefácio
\chapter*{Como usar esta apostila}
\addcontentsline{toc}{chapter}{Como usar esta apostila}
\markboth{Como usar esta apostila}{}

Esta apostila existe porque o HADROS3 usa espalhamento profundamente inelástico
(DIS, de \emph{deep inelastic scattering}) como um bloco fechado: uma seção de
choque $\sigma_{\nu N}(E)$ que sai de uma tabela e entra numa integral de
profundidade óptica. Isso é suficiente para \emph{rodar} o código, e insuficiente
para \emph{confiar} nele. O que está tabelado ali é o resultado final de sessenta
anos de física, e quem vai usar o HADROS3 precisa saber o que aquele número
significa, de onde ele vem e onde ele é frágil.

\medskip
\noindent\textbf{O que assumimos que você já sabe.} Relatividade restrita com
quadrivetores, mecânica quântica em nível de pós-graduação, e o mínimo de teoria
quântica de campos para reconhecer um espinor de Dirac e uma matriz $\gamma^\mu$.
Nada além disso: os traços de Dirac de que precisamos estão feitos passo a passo
no \cref{cap:tensores}.

\medskip
\noindent\textbf{Como o texto é organizado.} As deduções longas ficam em caixas
verdes; elas podem ser puladas numa primeira leitura, mas foram escritas
justamente para \emph{não} serem puladas na segunda. As caixas laranja marcam
armadilhas --- lugares onde é fácil escrever a fórmula certa e usá-la errado. As
caixas azuis são definições que serão usadas o tempo todo depois.

\medskip
\noindent\textbf{Verificações.} Sempre que possível, uma afirmação vem
acompanhada de um jeito de conferi-la: um limite conhecido, uma regra de soma,
ou um número que você pode recalcular. Duas das figuras deste texto são
construídas a partir dos dados públicos combinados de H1 e ZEUS~\cite{H1:2015ubc};
uma terceira usa as distribuições de pártons HERAPDF2.0 extraídas da mesma
publicação. O código que baixa esses dados, extrai os subconjuntos e desenha as
figuras está em \texttt{sandbox/apostila/figuras/}.

\medskip
\noindent\textbf{Onde isto se conecta ao HADROS3.} O \cref{cap:cc} fecha o
círculo: a fórmula de DIS de corrente carregada, a seção de choque de neutrinos
de energia ultra-alta, e a integral de profundidade óptica que o notebook
\texttt{01\_profundidade\_optica\_dis.ipynb} calcula. Recomendamos ler a apostila
até o \cref{cap:partons} e então fazer o notebook, voltando depois para os
\cref{cap:qcd,cap:cc}.

\tableofcontents

\cleardoublepage
\pagenumbering{arabic}

% =====================================================================
\chapter{Por que ``profundamente inelástico''?}
\label{cap:porque}

\section{A ideia é velha: jogue coisas e veja o que volta}

Em 1911 Geiger e Marsden jogaram partículas $\alpha$ contra uma folha de ouro.
A esmagadora maioria passou quase reto; uma fração minúscula voltou para trás.
Rutherford entendeu que o espalhamento a grande ângulo só era possível se a
carga positiva do átomo estivesse concentrada num caroço muito pequeno. O
método --- \emph{sonde com um projétil bem entendido e leia a estrutura do alvo
na distribuição angular} --- é o mesmo até hoje.

A diferença é a régua. Um projétil de momento transferido $q$ resolve estruturas
de tamanho
\begin{equation}
  \lambda \sim \frac{\hbar c}{\sqrt{\Qsq}},
  \qquad \Qsq \equiv -q^2 > 0 .
  \label{eq:regua}
\end{equation}
Com $\hbar c = 0{,}1973\gev\cdot\mathrm{fm}$, sondar o próton inteiro
($\approx 0{,}84$ fm) exige $\sqrt{\Qsq}\sim 0{,}2\gev$; enxergar dez vezes mais
fundo exige cem vezes mais $\Qsq$. É por isso que a física de estrutura é uma
física de aceleradores cada vez maiores.%
\nota{$1\,$fm $=10^{-15}$\,m. O próton tem raio de carga $0{,}841$\,fm; um núcleo de ferro, cerca de $4$\,fm.}

\nota{A \cref{fig:xtipico} mostra as duas escalas juntas: o $x$ sondado e o
$\lambda$ correspondente.}

\begin{figure}[t]
  \centering
  \fig{fig_x_tipico.pdf}
  \caption{À esquerda, a fração de momento $x$ que um neutrino de energia $E_\nu$
  tipicamente sonda; à direita, a resolução espacial associada a $\Qsq$. Um
  neutrino de $10^{9}\gev$ enxerga pártons com $x\sim 10^{-6}$ --- três ordens de
  grandeza abaixo do menor $x$ jamais medido.}
  \label{fig:xtipico}
\end{figure}

\section{Quatro regimes, um mesmo processo}

Considere um lépton batendo num núcleon. Chame de $W$ a massa invariante de
tudo o que sai do lado hadrônico. Conforme $\Qsq$ e $W$ crescem, o processo passa
por quatro regimes qualitativamente distintos:

\begin{description}[style=nextline,leftmargin=0pt,font=\sffamily\bfseries\color{acento}]
\item[Elástico ($W = M$)] O núcleon sobrevive inteiro. A informação está nos
  fatores de forma $G_E(\Qsq)$, $G_M(\Qsq)$, que são as transformadas de Fourier
  das distribuições de carga e magnetização. Foi assim que Hofstadter mediu o
  raio do próton nos anos 1950.
\item[Ressonante ($M < W \lesssim 2\gev$)] O núcleon é excitado para estados
  ligados --- $\Delta(1232)$ e companhia. O espectro é cheio de picos.
\item[Profundamente inelástico ($W \gg M$ \emph{e} $\Qsq \gg M^2$)] O núcleon é
  estilhaçado. Os picos somem e sobra um contínuo suave. É aqui que vivemos.
\item[Difrativo / $x$ pequeno] Um regime dentro do DIS em que a troca carrega os
  números quânticos do vácuo e o próton pode até sobreviver. Voltaremos a ele no
  \cref{cap:qcd}, porque é exatamente o regime dos neutrinos UHE.
\end{description}

O adjetivo ``profundamente'' se refere a $\Qsq \gg M^2$: a régua é muito menor
que o alvo. O adjetivo ``inelástico'' se refere a $W \gg M$: o alvo não sobrevive.
As duas condições juntas são o que permite a simplificação decisiva do
\cref{cap:partons} --- tratar o núcleon como uma \emph{coleção de constituintes
livres} durante o instante da colisão.

\begin{historia}[O experimento que mudou tudo]
Em 1968--69, no SLAC, o espectrômetro de 20\,GeV mediu espalhamento inelástico
$ep$ a $6^\circ$ e $10^\circ$~\cite{Bloom:1969kc,Breidenbach:1969kd}. Esperava-se
que a seção de choque caísse rapidamente com $\Qsq$, como cai no caso elástico,
porque um objeto extenso e mole não deveria conseguir absorver muito momento sem
se desmanchar suavemente. O que se viu foi o contrário: a seção de choque
inelástica caía muito devagar, e a razão em relação ao espalhamento por um alvo
puntual era praticamente independente de $\Qsq$.

Bjorken já havia previsto essa independência --- o \emph{escalonamento} --- a
partir de álgebra de correntes~\cite{Bjorken:1968dy}, e Feynman deu a ela a
imagem que usamos até hoje: o próton é feito de constituintes puntuais que, no
instante da colisão, se comportam como livres~\cite{Feynman:1969ej}. Friedman,
Kendall e Taylor receberam o Nobel de 1990 por esses
experimentos~\cite{Friedman:1972sy}.
\end{historia}

\section{O que realmente se mede}

Vale insistir num ponto que costuma escapar. Num experimento de DIS inclusivo
mede-se \emph{apenas o lépton que sai}: sua energia e seu ângulo. Nada do estado
hadrônico $X$ é observado --- soma-se sobre tudo o que ele possa ser. Isso tem
duas consequências:

\begin{enumerate}[leftmargin=*]
\item Toda a informação cabe em duas variáveis (a energia e o ângulo do lépton,
  ou equivalentemente $\Qsq$ e $\nu$). É por isso que a seção de choque é
  \emph{duplo}-diferencial e não mais que isso.
\item A soma sobre $X$ é o que torna possível escrever o lado hadrônico em
  termos de poucas funções --- as \emph{funções de estrutura}. Sem essa soma,
  não haveria simplificação alguma.
\end{enumerate}

% =====================================================================
\chapter{Cinemática}
\label{cap:cinematica}

Este capítulo é puro Lorentz: nenhuma dinâmica, nenhuma hipótese sobre
constituintes. Tudo aqui vale exatamente, sempre.

\section{O processo e seus quadrimomentos}

\begin{figure}[t]
\centering
\begin{tikzpicture}[
  x=1cm, y=1cm, every node/.style={font=\small},
  fermion/.style={-{Stealth[length=2.6mm]},thick,tinta},
  boson/.style={decorate,decoration={snake,amplitude=0.7mm,segment length=2.4mm},
                thick,acento},
  blob/.style={ellipse,draw=tinta,fill=acento!14,minimum width=14mm,
               minimum height=9mm,thick}
]
\coordinate (V) at (3.4,2.6);          % vértice leptônico
\coordinate (B) at (3.4,0.95);         % onde o bóson encontra o hádron

% --- lado leptônico ---------------------------------------------------
\draw[fermion] (0.35,3.35) -- node[above=1pt,pos=0.55]{$k$} (V);
\draw[fermion] (V) -- node[above=1pt,pos=0.45]{$k'$} (6.55,3.5);
\node[left]  at (0.3,3.35) {$\ell$};
\node[right] at (6.6,3.5)  {$\ell'$};
% continuação tracejada da direção incidente, para marcar o ângulo
\draw[suave,thin,dashed] (V) -- ++(-17.4:1.7);
\draw[suave,thin] ([shift=(-17.4:1.15)]V) arc (-17.4:15.7:1.15);
\node[suave,font=\footnotesize] at ([shift=(-1:1.42)]V) {$\theta$};

% --- bóson ------------------------------------------------------------
\draw[boson] (V) -- (B);
\node[anchor=east,font=\footnotesize] at (3.22,1.78) {$q = k-k'$};
\node[anchor=west,acento,font=\footnotesize] at (3.58,1.78) {$\gamma^*,\ Z,\ W^\pm$};

% --- lado hadrônico ---------------------------------------------------
\node[blob] (b) at (3.4,0.55) {};
\draw[fermion] (0.35,0.55) -- node[above=1pt]{$P$} (b.west);
\node[left] at (0.3,0.55) {$N$};
\foreach \a in {-30,-15,0,15,30}{
  \draw[thick,tinta,opacity=0.72] (b.east) -- ++(\a:2.05);}
\node[anchor=west] at (6.32,0.55) {$X$};
\node[suave,font=\footnotesize,anchor=north] at (4.5,-0.95)
  {$W^2 = (P+q)^2$};
\end{tikzpicture}
\caption{DIS inclusivo. Mede-se apenas $k'$; soma-se sobre todos os estados
finais hadrônicos $X$. A troca é um fóton virtual (corrente neutra
eletromagnética), um $Z$ (corrente neutra fraca) ou um $W^\pm$ (corrente
carregada). Esta apostila trata sobretudo dos casos $\gamma^*$ e $W^\pm$.}
\label{fig:processo}
\end{figure}

Fixemos a notação, que seguirá o resto do texto:
\[
\begin{array}{llll}
  k & \text{lépton incidente} &\qquad P & \text{núcleon alvo, } P^2 = M^2\\
  k' & \text{lépton espalhado} &\qquad q = k-k' & \text{momento transferido}
\end{array}
\]
Trabalhamos sempre com $\hbar = c = 1$ e desprezamos as massas leptônicas
(justificaremos isso no \cref{cap:cc}).

\section{Os invariantes}

\begin{definicao}[As variáveis de DIS]
\begin{align}
  s      &= (k+P)^2 \simeq 2\,M E_\ell &&\text{energia de centro de massa ao quadrado}\\
  \Qsq   &= -q^2 \;=\; 4EE'\sin^2(\theta/2) &&\text{virtualidade da troca}\\
  \nu    &= \frac{P\cdot q}{M} \;=\; E - E' &&\text{energia transferida (referencial do alvo)}\\
  \xB    &= \frac{\Qsq}{2P\cdot q} = \frac{\Qsq}{2M\nu} &&\text{variável de Bjorken}\\
  y      &= \frac{P\cdot q}{P\cdot k} = \frac{\nu}{E} &&\text{inelasticidade}\\
  W^2    &= (P+q)^2 &&\text{massa invariante do sistema hadrônico}
\end{align}
As formas $\Qsq = 4EE'\sin^2(\theta/2)$, $\nu = E-E'$ e $y=\nu/E$ valem no
referencial de repouso do alvo; $s$, $\Qsq$, $\xB$, $y$ e $W^2$ são invariantes
de Lorentz e podem ser usados em qualquer referencial. Guarde essa distinção: ela
volta a importar quando o ``referencial do alvo'' passa a ser o de um fluido
girando perto de um buraco negro, como no HADROS3.%
\nota{No HADROS3 isso vira a energia local $E=-u^\mu p_\mu$, medida por um observador estático ou ZAMO. Ver \texttt{dis\_sampler.py:185--203}.}
\end{definicao}

Só duas dessas variáveis são independentes (dado $s$). As relações entre elas
são o assunto do resto da seção.

\subsection{A relação $\Qsq = xys$}

\begin{deducao}[Dedução 2.1]
Por definição, $\xB = \Qsq/(2P\cdot q)$, logo
\begin{equation}
  \Qsq = 2\,\xB\,(P\cdot q).
\end{equation}
Também por definição $y = (P\cdot q)/(P\cdot k)$, então $P\cdot q = y\,(P\cdot k)$:
\begin{equation}
  \Qsq = 2\,\xB\,y\,(P\cdot k).
\end{equation}
Falta identificar $P\cdot k$. Como $s = (k+P)^2 = k^2 + 2P\cdot k + P^2$ e estamos
desprezando a massa do lépton ($k^2=0$),
\begin{equation}
  s = 2P\cdot k + M^2 \quad\Longrightarrow\quad P\cdot k = \frac{s-M^2}{2}.
\end{equation}
Portanto, \emph{exatamente},
\begin{equation}
  \boxed{\ \Qsq = \xB\,y\,(s-M^2)\ } \qquad
  \xrightarrow{\ s\,\gg\, M^2\ }\qquad \Qsq \simeq \xB\,y\,s .
  \label{eq:q2xys}
\end{equation}
\end{deducao}

\noindent Essa é a relação mais usada de todo o assunto, e ela diz uma coisa
importante: \emph{fixado $s$, escolher $(x,y)$ é escolher $\Qsq$}. As três
variáveis não são independentes.

\subsection{Por que $0 \le x \le 1$}

\begin{deducao}[Dedução 2.2]
Escreva a massa invariante do sistema hadrônico:
\begin{equation}
  W^2 = (P+q)^2 = P^2 + 2P\cdot q + q^2 = M^2 + 2M\nu - \Qsq .
  \label{eq:W2}
\end{equation}
Agora use $2M\nu = \Qsq/\xB$, que é a definição de $\xB$ reescrita:
\begin{equation}
  W^2 = M^2 + \frac{\Qsq}{\xB} - \Qsq
      = M^2 + \Qsq\,\frac{1-\xB}{\xB}.
  \label{eq:W2x}
\end{equation}
O estado final hadrônico mais leve possível é o próprio núcleon, então
$W^2 \ge M^2$. Com $\Qsq>0$, a \cref{eq:W2x} exige
\begin{equation}
  \frac{1-\xB}{\xB}\ \ge\ 0 \quad\Longrightarrow\quad \boxed{\ 0 < \xB \le 1\ }.
\end{equation}
E o extremo é informativo: $\xB=1$ dá exatamente $W^2=M^2$, isto é,
\textbf{$x=1$ é o espalhamento elástico}. O regime profundamente inelástico é
$x$ bem abaixo de 1 com $\Qsq$ grande.%
\nota{Limites úteis para conferir contas: $x=1$ é elástico; $y=0$ é o lépton passando reto; $\Qsq=0$ é um fóton real.}
\end{deducao}

\begin{cuidado}[Duas coisas diferentes com o mesmo nome]
A variável $\xB$ definida acima é puramente cinemática --- ela existe mesmo que o
próton não tenha constituinte nenhum. No \cref{cap:partons} vamos mostrar que,
\emph{sob as hipóteses do modelo de pártons}, $\xB$ coincide com a fração de
momento do párton atingido. Coincidem, mas não são a mesma coisa: a primeira é
uma definição, a segunda é um resultado físico com condições de validade. Em
correções de ordem superior, e sempre que massas importam, elas se separam.
\end{cuidado}

\subsection{O intervalo de $y$ e o significado físico}

No referencial de repouso do alvo, $y = \nu/E = (E-E')/E$ é a \emph{fração da
energia do lépton transferida ao sistema hadrônico}. Como $0 \le E' \le E$, temos
imediatamente $0 \le y \le 1$. Em $y\to0$ o lépton mal é perturbado; em $y\to1$
ele entrega tudo.

Há ainda um vínculo conjunto: como $\Qsq = xys \le s$, nem todo par $(x,y)$ é
acessível a um dado $s$. O lugar geométrico dessa restrição é o que a
\cref{fig:plano} mostra.

\begin{figure}[t]
  \centering
  \fig{fig_plano_cinematico.pdf}
  \caption{O plano $(x,\Qsq)$. Os pontos são as $485$ medidas combinadas de H1 e
  ZEUS para corrente neutra $e^+p$~\cite{H1:2015ubc}; a faixa clara é a região
  típica de experimentos de alvo fixo. As linhas pontilhadas marcam o $x$ que um
  neutrino de energia $E_\nu$ sonda quando o propagador do $W$ fixa
  $\Qsq\sim m_W^2$. A distância entre os dados e essas linhas é o problema
  central da física de neutrinos UHE.}
  \label{fig:plano}
\end{figure}

\section{Trocando de variáveis: o jacobiano}

Experimentalmente é natural medir $(\Qsq,\nu)$ --- ou $(E',\theta)$. Teoricamente
é natural usar $(x,y)$. A conversão aparece o tempo todo e vale fazer com cuidado.

\begin{deducao}[Dedução 2.3]
No referencial de repouso do alvo, com $E$ fixo,
\begin{equation}
  \Qsq = 2ME\,x\,y, \qquad \nu = E\,y .
\end{equation}
O jacobiano da transformação $(x,y)\mapsto(\Qsq,\nu)$ é o determinante
\begin{equation}
  \left|\frac{\partial(\Qsq,\nu)}{\partial(x,y)}\right|
  = \left|\det\!\begin{pmatrix}
      \partial \Qsq/\partial x & \partial \Qsq/\partial y\\[2pt]
      \partial \nu /\partial x & \partial \nu /\partial y
    \end{pmatrix}\right|
  = \left|\det\!\begin{pmatrix} 2MEy & 2MEx\\ 0 & E \end{pmatrix}\right|
  = 2ME^{2}y .
\end{equation}
Logo
\begin{equation}
  \frac{\md^2\sigma}{\md x\,\md y} = 2ME^{2}y\;\frac{\md^2\sigma}{\md \Qsq\,\md \nu}.
  \label{eq:jacobiano}
\end{equation}
\end{deducao}

\begin{cuidado}[O fator $y$ some se você não prestar atenção]
O jacobiano tem um $y$ dentro. Se você comparar uma distribuição em $y$ obtida de
$\md^2\sigma/\md x\,\md y$ com uma obtida de $\md^2\sigma/\md \Qsq\,\md\nu$ sem
incluí-lo, vai concluir que a física está errada quando na verdade a conta é que
está. Este é, empiricamente, o erro mais comum de quem começa.
\end{cuidado}

% =====================================================================
\chapter{A seção de choque em forma tensorial}
\label{cap:tensores}

Este é o capítulo técnico. O objetivo é chegar em
$\md^2\sigma/\md x\,\md y$ escrita em termos de funções de estrutura, sem
nenhuma hipótese sobre o que há dentro do núcleon. A estratégia é a mesma em
qualquer livro-texto~\cite{Halzen:1984mc,Devenish:2004pb,Ellis:1996mzs}: separar
o que sabemos calcular (o lado leptônico) do que não sabemos (o lado hadrônico),
e espremer o segundo até que só sobre o que a simetria permite.

\section{A fatoração}

Em ordem mais baixa há um único diagrama: o lépton emite um bóson virtual, que é
absorvido pelo núcleon. A amplitude fatoriza,
\begin{equation}
  \mathcal{M} \;=\; \underbrace{\bar u(k')\,\Gamma^\mu\,u(k)}_{\text{lado leptônico}}
  \;\times\; \underbrace{\frac{-i\,g_{\mu\nu}}{q^2 - m_B^2}}_{\text{propagador}}
  \;\times\; \underbrace{\langle X|J^\nu|P\rangle}_{\text{lado hadrônico}},
\end{equation}
e ao tomar $|\mathcal{M}|^2$, somar sobre estados finais e sobre spins, a seção
de choque fica proporcional a uma contração de dois tensores:
\begin{equation}
  \boxed{\ \frac{\md^2\sigma}{\md \Qsq\,\md\nu} \;\propto\; L^{\mu\nu}\,W_{\mu\nu}\ }
  \label{eq:LW}
\end{equation}

O tensor leptônico $L^{\mu\nu}$ é calculável exatamente --- é
eletrodinâmica/teoria eletrofraca pura.%
\nota{A separação em dois tensores só é possível porque trocamos \emph{um} bóson. Em ordens mais altas ela deixa de valer literalmente.} O tensor hadrônico $W_{\mu\nu}$ contém
toda a QCD não perturbativa e \emph{não} é calculável a partir de primeiros
princípios. A saída é não tentar calculá-lo: apenas descobrir quantas funções
independentes ele pode conter.

\section{O tensor leptônico, passo a passo}

\begin{deducao}[Dedução 3.1 --- caso eletromagnético]
Para um elétron, o vértice é $\Gamma^\mu = \gamma^\mu$ e
\begin{equation}
  L^{\mu\nu} = \frac{1}{2}\sum_{\text{spins}}
    \big[\bar u(k')\gamma^\mu u(k)\big]\big[\bar u(k')\gamma^\nu u(k)\big]^{*},
\end{equation}
onde o $\tfrac12$ é a média sobre os dois estados de spin do elétron incidente.
Usando $[\bar u(k')\gamma^\nu u(k)]^{*} = \bar u(k)\gamma^\nu u(k')$ e as relações
de completeza $\sum_s u(p)\bar u(p) = \sla{p} + m$, com $m\to0$:
\begin{equation}
  L^{\mu\nu} = \frac{1}{2}\,\mathrm{Tr}\!\left[\sla{k}'\gamma^\mu \sla{k}\gamma^\nu\right].
\end{equation}
Escrevendo $\sla{k}' = k'_\alpha\gamma^\alpha$ e $\sla{k} = k_\beta\gamma^\beta$,
precisamos do traço de quatro matrizes gama,
\begin{equation}
  \mathrm{Tr}\!\left[\gamma^\alpha\gamma^\mu\gamma^\beta\gamma^\nu\right]
  = 4\left(g^{\alpha\mu}g^{\beta\nu} - g^{\alpha\beta}g^{\mu\nu}
           + g^{\alpha\nu}g^{\mu\beta}\right).
  \label{eq:traco4}
\end{equation}
Contraindo com $k'_\alpha k_\beta$:%
\nota{A \cref{eq:traco4} é a identidade de traço mais usada da área. Vale a pena memorizá-la: dois $g$ com sinal $+$, o do meio com sinal $-$.}
\begin{equation}
  \mathrm{Tr}\!\left[\sla{k}'\gamma^\mu \sla{k}\gamma^\nu\right]
  = 4\left(k'^{\mu}k^{\nu} - (k'\!\cdot\!k)\,g^{\mu\nu} + k'^{\nu}k^{\mu}\right).
\end{equation}
Finalmente, com $\Qsq = -(k-k')^2 = 2\,k\!\cdot\!k'$ (massas nulas),
\begin{equation}
  \boxed{\;L^{\mu\nu}_{\rm em}
   = 2\left(k^{\mu}k'^{\nu} + k^{\nu}k'^{\mu}\right) - \Qsq\,g^{\mu\nu}\;}
  \label{eq:Lem}
\end{equation}
\end{deducao}

Repare em duas propriedades da \cref{eq:Lem} que vão ser usadas:
\begin{enumerate}[leftmargin=*,itemsep=2pt]
\item $L^{\mu\nu}$ é \emph{simétrico} em $\mu\leftrightarrow\nu$;
\item $q_\mu L^{\mu\nu} = 0$ quando as massas leptônicas são nulas. De fato,
  $q_\mu[2(k^\mu k'^\nu + k^\nu k'^\mu) - \Qsq g^{\mu\nu}]
   = 2[(q\!\cdot\!k)k'^\nu + (q\!\cdot\!k')k^\nu] - \Qsq q^\nu$, e como
  $q\cdot k = \Qsq/2$ e $q\cdot k' = -\Qsq/2$ (verifique!), isso é
  $\Qsq(k'^\nu - k^\nu) - \Qsq q^\nu = 0$.
\end{enumerate}

\begin{deducao}[Dedução 3.2 --- caso de corrente carregada]
Para um neutrino a estrutura do vértice é $V-A$: $\Gamma^\mu = \gamma^\mu(1-\gamma^5)$.
Não há média sobre spin --- o neutrino só existe com hélicidade esquerda. Então
\begin{equation}
  L^{\mu\nu}_{\rm CC}
  = \mathrm{Tr}\!\left[\sla{k}'\gamma^\mu(1-\gamma^5)\,\sla{k}\,\gamma^\nu(1-\gamma^5)\right].
\end{equation}
Use $(1-\gamma^5)\sla{k} = \sla{k}(1+\gamma^5)$ e $(1+\gamma^5)(1-\gamma^5)=0$,
$(1-\gamma^5)^2 = 2(1-\gamma^5)$, para escrever
\begin{equation}
  L^{\mu\nu}_{\rm CC}
  = 2\,\mathrm{Tr}\!\left[\sla{k}'\gamma^\mu \sla{k}\gamma^\nu(1-\gamma^5)\right].
\end{equation}
A parte sem $\gamma^5$ é o traço que já fizemos. A parte com $\gamma^5$ usa
\begin{equation}
  \mathrm{Tr}\!\left[\gamma^\alpha\gamma^\mu\gamma^\beta\gamma^\nu\gamma^5\right]
  = -4i\,\epsilon^{\alpha\mu\beta\nu},
  \qquad (\epsilon^{0123} = +1),
\end{equation}
e o resultado é
\begin{equation}
  \boxed{\;L^{\mu\nu}_{\rm CC}
   = 8\left[k^{\mu}k'^{\nu} + k^{\nu}k'^{\mu}
     - g^{\mu\nu}(k\!\cdot\!k')
     \;{\color{acento2}\mp}\; i\,\epsilon^{\mu\nu\alpha\beta}k_\alpha k'_\beta\right]\;}
  \label{eq:Lcc}
\end{equation}
com o sinal superior para neutrinos e o inferior para antineutrinos.
\end{deducao}

\begin{cuidado}[O termo antissimétrico é a paridade]
A diferença essencial entre \cref{eq:Lem} e \cref{eq:Lcc} é o termo com
$\epsilon^{\mu\nu\alpha\beta}$, que é \emph{antissimétrico} em
$\mu\leftrightarrow\nu$. Ele existe porque a interação fraca viola paridade, e é
ele que dará origem à terceira função de estrutura, $xF_3$ --- a que distingue
neutrino de antineutrino. Se você já sabe que $xF_3$ mede $q-\bar q$, essa é a
raiz: quarks e antiquarks acoplam de forma diferente a uma corrente que
distingue hélicidade.
\end{cuidado}

\section{O tensor hadrônico: o que a simetria permite}

Aqui está o truque central. Definimos
\begin{equation}
  W^{\mu\nu}(P,q) = \frac{1}{4\pi}\sum_X \int \md\Pi_X\,
    \langle P|J^{\mu\dagger}|X\rangle\langle X|J^{\nu}|P\rangle\,
    (2\pi)^4\delta^4(P+q-p_X),
\end{equation}
com a soma sobre \emph{todos} os estados hadrônicos --- é a soma inclusiva do
\cref{cap:porque}. Não sabemos calcular isso. Mas sabemos que $W^{\mu\nu}$ é um
tensor de rank 2 que só pode ser construído com os objetos disponíveis:%
\nota{Este é o padrão argumentativo mais poderoso da física de partículas: não calcule o que não sabe --- conte quantas funções a simetria permite, e meça-as.}
\begin{equation}
  g^{\mu\nu},\qquad P^\mu,\qquad q^\mu,\qquad \epsilon^{\mu\nu\alpha\beta}P_\alpha q_\beta .
\end{equation}
(Não há mais nada: o alvo não é polarizado, então não existe vetor de spin.) A
forma mais geral compatível com invariância de Lorentz é, portanto,
\begin{align}
  W^{\mu\nu} =\;& -W_1\,g^{\mu\nu}
   + \frac{W_2}{M^2}\,P^\mu P^\nu
   - \frac{i\,W_3}{2M^2}\,\epsilon^{\mu\nu\alpha\beta}P_\alpha q_\beta \nonumber\\
   &+ \frac{W_4}{M^2}\,q^\mu q^\nu
   + \frac{W_5}{2M^2}\left(P^\mu q^\nu + q^\mu P^\nu\right)
   + \frac{i\,W_6}{2M^2}\left(P^\mu q^\nu - q^\mu P^\nu\right),
  \label{eq:Wgeral}
\end{align}
onde cada $W_i = W_i(\Qsq,\nu)$ é uma função escalar --- duas variáveis, porque
$P^2 = M^2$ é fixo e só sobram $q^2$ e $P\cdot q$ como invariantes.

\begin{deducao}[Dedução 3.3 --- o corte de seis para três]
Duas observações eliminam metade dos termos.

\smallskip
\textbf{(i) Conservação da corrente.} A corrente eletromagnética é conservada,
$\partial_\mu J^\mu = 0$, o que impõe
\begin{equation}
  q_\mu W^{\mu\nu} = 0 = q_\nu W^{\mu\nu}.
\end{equation}
Aplicando isso à \cref{eq:Wgeral} obtêm-se relações que expressam $W_4$, $W_5$ e
$W_6$ em termos de $W_1$ e $W_2$. (Para a corrente carregada a corrente não é
exatamente conservada, e $W_4$, $W_5$ sobrevivem --- mas seus coeficientes na
contração final vêm multiplicados por $m_\ell^2$, a massa do lépton carregado
produzido. Para $E_\nu$ de GeV para cima isso é desprezível, exceto para o
$\tau$.)

\smallskip
\textbf{(ii) O que sobrevive à contração.} Ainda que $W_4$, $W_5$, $W_6$ não
fossem elimináveis, seus termos são proporcionais a $q^\mu$ ou $q^\nu$. Como
mostramos que $q_\mu L^{\mu\nu} = 0$ no limite de lépton sem massa, esses termos
\emph{não contribuem para nada de observável}.

\smallskip
Sobram três: $W_1$, $W_2$, $W_3$. Este é o conteúdo do teorema --- o lado
hadrônico do DIS inclusivo não polarizado tem exatamente três graus de liberdade,
não importa quão complicada seja a QCD lá dentro.
\end{deducao}

\begin{definicao}[Funções de estrutura adimensionais]
É costume trocar $W_{1,2,3}$ por combinações adimensionais:
\begin{equation}
  F_1(x,\Qsq) = W_1,\qquad
  F_2(x,\Qsq) = \frac{\nu}{M}\,W_2,\qquad
  F_3(x,\Qsq) = \frac{\nu}{M}\,W_3 .
\end{equation}
A vantagem aparece no \cref{cap:partons}: nessas variáveis, o modelo de pártons
prevê que as $F_i$ dependem \emph{só de $x$}, e não de $\Qsq$.
\end{definicao}

\section{A contração e a fórmula mestra}

Contraindo \cref{eq:Lcc} com \cref{eq:Wgeral} e convertendo para as variáveis
$(x,y)$ com o jacobiano da \cref{eq:jacobiano}, chega-se ao resultado central de
toda a área:

\begin{tcolorbox}[caixabase=acento,title={A fórmula mestra do DIS de corrente carregada}]
\begin{equation}
  \frac{\md^2\sigma^{\rm CC}_{\nu(\bar\nu)N}}{\md x\,\md y}
  = \frac{G_F^2\,s}{4\pi}
    \left(\frac{m_W^2}{m_W^2 + \Qsq}\right)^{\!2}
    \Big[\,Y_+\,F_2(x,\Qsq)\;-\;y^2 F_L(x,\Qsq)\;{\pm}\;Y_-\,xF_3(x,\Qsq)\Big],
  \label{eq:mestra}
\end{equation}
com
\begin{equation}
  Y_\pm \equiv 1 \pm (1-y)^2,
  \qquad
  F_L \equiv F_2 - 2xF_1 ,
\end{equation}
sinal $+$ para $\nu$ e $-$ para $\bar\nu$. Para o caso eletromagnético
($e^\pm N$) troque $G_F^2 s/4\pi \to 4\pi\alpha^2 s/(\Qsq)^{2}\cdot(\dots)$,
elimine $F_3$ e o fator de propagador.
\end{tcolorbox}

\begin{deducao}[Dedução 3.4 --- por que esta forma é igual à do PDG]
O PDG~\cite{ParticleDataGroup:2024cfk} escreve a mesma coisa assim:
\begin{equation}
  \frac{\md^2\sigma}{\md x\,\md y}
  = \frac{G_F^2 s}{2\pi}\left(\frac{m_W^2}{m_W^2+\Qsq}\right)^{\!2}
    \Big[(1-y)F_2 + y^2 xF_1 \pm y\big(1-\tfrac{y}{2}\big)xF_3\Big].
  \label{eq:pdgform}
\end{equation}
Parecem diferentes; são idênticas. Substituindo $F_L = F_2 - 2xF_1$ no colchete
da \cref{eq:mestra}:
\begin{align}
  Y_+F_2 - y^2F_L
  &= \big[1+(1-y)^2\big]F_2 - y^2\big(F_2 - 2xF_1\big)\nonumber\\
  &= F_2\big[\underbrace{1+(1-y)^2-y^2}_{=\;2(1-y)}\big] + 2y^2\,xF_1
   = 2(1-y)F_2 + 2y^2 xF_1 .
\end{align}
E para o termo de $F_3$:
\begin{equation}
  Y_- = 1-(1-y)^2 = 2y - y^2 = 2y\big(1-\tfrac{y}{2}\big).
\end{equation}
Pondo o fator 2 comum em evidência, $\tfrac{G_F^2s}{4\pi}\cdot 2 = \tfrac{G_F^2s}{2\pi}$,
recupera-se exatamente a \cref{eq:pdgform}. As duas formas convivem na
literatura; a primeira é mais usada em HERA (porque isola $F_L$), a segunda em
física de neutrinos. É a mesma equação.
\end{deducao}

\begin{figure}[t]
  \centering
  \fig{fig_dependencia_y.pdf}
  \caption{Os pesos cinemáticos da \cref{eq:mestra}. À direita, $Y_+$ multiplica
  $F_2$, $Y_-$ multiplica $xF_3$ e $y^2$ multiplica $F_L$: em $y$ pequeno só
  $F_2$ sobrevive, e é por isso que a seção de choque reduzida medida a baixo $y$
  é essencialmente $F_2$. À esquerda, a distribuição em $y$ prevista pelo modelo
  de pártons (\cref{cap:partons}) para quarks e antiquarks.}
  \label{fig:ydep}
\end{figure}

% =====================================================================
\chapter{O modelo de pártons}
\label{cap:partons}

Até aqui, nenhuma hipótese sobre o interior do núcleon. Agora fazemos uma, e
tudo desaba num resultado de uma linha.

\section{A hipótese}

\begin{definicao}[Modelo de pártons de Feynman--Bjorken]
Num referencial em que o núcleon tem momento muito grande (``momento infinito''),
o núcleon é uma coleção de constituintes puntuais --- os \emph{pártons} ---
que, durante o tempo brevíssimo da colisão, se comportam como \emph{livres} e
\emph{colineares} com o núcleon. O párton $i$ carrega uma fração $\xi$ do momento
do núcleon, distribuída segundo uma densidade $f_i(\xi)$.
\end{definicao}

\noindent A hipótese tem uma justificativa física: no referencial de momento
infinito, a dilatação temporal congela as interações internas do núcleon, cujo
tempo característico é $\sim 1/M$, enquanto a colisão dura $\sim 1/\sqrt{\Qsq}$.%
\nota{Esta é a única justificativa física do modelo de pártons. Se $\Qsq$ não for grande, ele simplesmente não vale.}
Se $\Qsq \gg M^2$ --- e é exatamente isso que ``profundamente'' significa --- a
colisão é rápida demais para que os pártons ``percebam'' uns aos outros.

\section{Espalhamento por um párton livre}

\begin{deducao}[Dedução 4.1 --- $x$ é a fração de momento]
Seja $p = \xi P$ o momento do párton atingido. Depois de absorver o bóson de
momento $q$, o párton continua sendo um párton, de massa desprezível:
\begin{equation}
  (p + q)^2 = 0 .
\end{equation}
Expandindo, com $p^2 = \xi^2M^2 \approx 0$:
\begin{equation}
  2\,p\cdot q + q^2 = 0
  \quad\Longrightarrow\quad
  2\,\xi\,(P\cdot q) = \Qsq
  \quad\Longrightarrow\quad
  \xi = \frac{\Qsq}{2P\cdot q} = \xB .
\end{equation}
Ou seja: \textbf{a variável cinemática de Bjorken $\xB$ é a fração de momento do
párton que foi atingido}. Este é o resultado que dá sentido físico a tudo o que
veio antes. Note as hipóteses usadas: párton sem massa, colinear, e livre no
estado final.
\end{deducao}

\section{A seção de choque e as funções de estrutura}

Para a corrente carregada, o subprocesso elementar é $\nu q \to \ell^- q'$ --- é
espalhamento de duas partículas puntuais, e a seção de choque é elementar:
\begin{equation}
  \frac{\md\hat\sigma}{\md y}(\nu q) = \frac{G_F^2\,\hat s}{\pi},
  \qquad
  \frac{\md\hat\sigma}{\md y}(\nu \bar q) = \frac{G_F^2\,\hat s}{\pi}\,(1-y)^2,
  \label{eq:subproc}
\end{equation}
com $\hat s = x s$ a energia de centro de massa do subprocesso.

\begin{deducao}[Dedução 4.2 --- de onde vem o $(1-y)^2$]
A interação fraca de corrente carregada só acopla partículas de hélicidade
esquerda e antipartículas de hélicidade direita. No centro de massa do
subprocesso:
\begin{itemize}[leftmargin=*,itemsep=2pt]
\item $\nu q$: o neutrino tem hélicidade $-\tfrac12$, o quark também. O momento
  angular total ao longo do eixo de colisão é $J_z = 0$. Um estado com $J=0$ é
  isotrópico --- não há direção preferida --- logo $\md\hat\sigma/\md\cos\theta^*$
  é constante. Como $y = (1-\cos\theta^*)/2$, isso dá
  $\md\hat\sigma/\md y = \text{const}$.
\item $\nu\bar q$: agora o antiquark tem hélicidade $+\tfrac12$ e
  $J_z = -1$. O espalhamento para trás ($\theta^*=\pi$, isto é $y=1$) exigiria
  inverter $J_z$, o que é proibido. A amplitude carrega a função de Wigner
  $d^{1}_{-1,-1}(\theta^*) = \tfrac{1}{2}(1+\cos\theta^*) = 1-y$, e a seção de
  choque, o quadrado: $(1-y)^2$.
\end{itemize}
Essa é a assinatura mais limpa da estrutura $V-A$, e é o painel esquerdo da
\cref{fig:ydep}. Medir a distribuição em $y$ é medir diretamente o conteúdo de
antiquarks do núcleon.
\end{deducao}

Somando incoerentemente sobre os pártons, ponderados por suas densidades:
\begin{equation}
  \frac{\md^2\sigma^{\rm CC}_{\nu N}}{\md x\,\md y}
  = \frac{G_F^2\,x\,s}{\pi}\Big[\,q(x) + \bar q(x)\,(1-y)^2\,\Big].
  \label{eq:partonxsec}
\end{equation}

\begin{deducao}[Dedução 4.3 --- lendo as funções de estrutura]
Compare \cref{eq:partonxsec} com a forma do PDG, \cref{eq:pdgform} (sem o fator
de propagador, que é 1 em $\Qsq \ll m_W^2$). Escrevendo o \emph{ansatz}
\begin{equation}
  F_2 = 2x\big[q+\bar q\big], \qquad
  xF_3 = 2x\big[q-\bar q\big], \qquad
  2xF_1 = F_2 ,
\end{equation}
o colchete da \cref{eq:pdgform} fica
\begin{align}
  &(1-y)\,2x(q+\bar q) + y^2\,x\!\cdot\!\frac{2x(q+\bar q)}{2x}
    + y\Big(1-\frac{y}{2}\Big)2x(q-\bar q) \nonumber\\
  &\quad= 2x\Big\{ q\underbrace{\Big[(1-y) + \tfrac{y^2}{2} + y - \tfrac{y^2}{2}\Big]}_{=\,1}
   + \bar q\underbrace{\Big[(1-y) + \tfrac{y^2}{2} - y + \tfrac{y^2}{2}\Big]}_{=\,(1-y)^2}\Big\}.
\end{align}
Multiplicando pelo prefator $G_F^2s/2\pi$ obtém-se exatamente a
\cref{eq:partonxsec}. O \emph{ansatz} estava certo:
\begin{equation}
  \boxed{\;F_2(x) = 2x\sum_i \big[q_i(x) + \bar q_i(x)\big],\qquad
          xF_3(x) = 2x\sum_i \big[q_i(x) - \bar q_i(x)\big].\;}
\end{equation}
Para o caso eletromagnético, o mesmo raciocínio com o acoplamento $e_q$ em vez de
$G_F$ dá
\begin{equation}
  F_2^{\rm em}(x) = \sum_i e_i^2\, x\big[q_i(x) + \bar q_i(x)\big].
  \label{eq:F2em}
\end{equation}
\end{deducao}

\section{Escalonamento de Bjorken}

Olhe para o que acabou de acontecer: as funções de estrutura ficaram funções
\emph{só de $x$}. Nenhum $\Qsq$ aparece. Essa é a previsão de
Bjorken~\cite{Bjorken:1968dy,Bjorken:1969ja}, e a razão física é simples: se os
pártons são puntuais, não há escala de comprimento com que $\Qsq$ possa ser
comparado, e portanto $\Qsq$ não pode aparecer numa quantidade adimensional.
Escalonamento é a assinatura de constituintes puntuais, do mesmo modo que o
espalhamento de Rutherford é a assinatura de um núcleo puntual.%
\nota{Compare: no espalhamento elástico, o fator de forma cai rápido com $\Qsq$ porque o alvo é extenso. Aqui não cai --- porque o alvo relevante é puntual.}

\section{A relação de Callan--Gross}

\begin{deducao}[Dedução 4.4 --- o spin do párton está em $F_L$]
Defina a função de estrutura longitudinal $F_L = F_2 - 2xF_1$. Ela é proporcional
à seção de choque de absorção de um bóson virtual \emph{longitudinalmente}
polarizado.

Considere um párton sem massa de spin $\tfrac12$ absorvendo um fóton virtual.
A helicidade de um férmion sem massa é conservada pelo vértice $\gamma^\mu$.
Um fóton longitudinal tem helicidade $0$; absorvê-lo mudaria a projeção de spin
do quark de $\pm\tfrac12$ para $\pm\tfrac12$ com $\Delta J_z = 0$ ao longo do
eixo --- o que só seria compatível se o quark pudesse virar de helicidade, e ele
não pode. Logo
\begin{equation}
  \sigma_L = 0 \quad\Longrightarrow\quad F_L = 0
  \quad\Longrightarrow\quad \boxed{\;F_2(x) = 2x F_1(x)\;}
  \label{eq:cg}
\end{equation}
que é a relação de Callan--Gross~\cite{Callan:1969uq}.

O contraste é o que torna isso um \emph{teste}: se os pártons tivessem spin $0$,
não haveria acoplamento a fótons transversais, e o resultado seria o oposto:
$F_1 = 0$ e $F_L = F_2$. A medida experimental de $R = \sigma_L/\sigma_T \approx 0$
nos anos 1970 foi a evidência direta de que os pártons são férmions de spin
$\tfrac12$%
\nota{Um raro caso em que uma medida \emph{nula} é a descoberta.} --- ou seja, de que são os quarks de Gell-Mann e Zweig, e não algum
outro objeto.
\end{deducao}

\begin{cuidado}[$F_L$ não é zero de verdade]
Em QCD, $F_L$ é pequeno mas \emph{não nulo}: um quark pode emitir um glúon e
adquirir momento transverso, quebrando o argumento de helicidade. $F_L$ é, na
verdade, uma medida bastante direta da densidade de glúons, e foi medido em HERA.
Na \cref{eq:mestra} ele entra com peso $y^2$; é por isso que as análises que
querem $F_2$ limpo cortam em $y$ baixo --- como fizemos na \cref{fig:violacao}.
\end{cuidado}

\section{Regras de soma, e uma verificação numérica}

O modelo de pártons vem com vínculos que não dependem de nenhum ajuste. Se
$u_v \equiv u - \bar u$ e $d_v \equiv d - \bar d$ são as distribuições de
valência, contar quarks no próton exige
\begin{equation}
  \int_0^1 u_v(x)\,\md x = 2, \qquad \int_0^1 d_v(x)\,\md x = 1 .
  \label{eq:contagem}
\end{equation}
E somar o momento de todos os constituintes tem de dar o momento do próton:
\begin{equation}
  \int_0^1 x\Big[\sum_i q_i(x) + \bar q_i(x) + g(x)\Big]\md x = 1 .
  \label{eq:momento}
\end{equation}

\begin{verifique}[Confira você mesmo]
As regras de contagem da \cref{eq:contagem} não são aproximações --- são
identidades que qualquer conjunto de PDFs tem de satisfazer. Usando o conjunto
HERAPDF2.0 LO~\cite{H1:2015ubc}, extraído em
\texttt{sandbox/data/hera/herapdf20\_lo\_xf.dat}, uma integração numérica direta
dá
\[
  \int_0^1 u_v\,\md x = 2{,}0008, \qquad
  \int_0^1 d_v\,\md x = 1{,}0003 \qquad (\Qsq = 10\gev^2),
\]
e os mesmos números até a terceira casa em $\Qsq = 10^4\gev^2$ --- as regras de
soma são preservadas pela evolução em $\Qsq$, como têm de ser. O desvio de
$4\times10^{-4}$ é do integrador, não da física.

Historicamente, a \cref{eq:momento} deu um susto: a soma sobre os \emph{quarks}
dá cerca de $0{,}5$, não $1$. Metade do momento do próton é carregada por algo
que não interage eletromagneticamente nem fracamente. Esse ``algo'' é o glúon.
\end{verifique}

\begin{figure}[t]
  \centering
  \fig{fig_pdfs.pdf}
  \caption{Distribuições de pártons do próton, HERAPDF2.0 LO~\cite{H1:2015ubc},
  em duas escalas. Note o que a evolução faz: em $\Qsq$ maior, a valência
  encolhe e migra para $x$ pequeno, enquanto o mar e o glúon explodem. O glúon
  está dividido por 10 para caber no gráfico.}
  \label{fig:pdfs}
\end{figure}

% =====================================================================
\chapter{QCD: as distribuições evoluem}
\label{cap:qcd}

\section{O escalonamento quebra --- e isso é bom}

O escalonamento de Bjorken vale se os pártons forem realmente puntuais e
realmente livres. Nenhuma das duas coisas é exatamente verdade: os quarks emitem
glúons. E emitir um glúon quase colinear é quase grátis, o que introduz uma
dependência --- logarítmica --- em $\Qsq$.

\begin{deducao}[Dedução 5.1 --- de onde vem o logaritmo]
Considere um quark de momento $p$ emitindo um glúon e ficando com fração $z$.
A integral sobre o momento transverso $k_T$ do glúon emitido tem a forma
\begin{equation}
  \int \frac{\md^2 k_T}{k_T^2}\;\sim\;\pi\int^{\Qsq}_{\mu^2}\frac{\md k_T^2}{k_T^2}
  \;=\;\pi\,\ln\frac{\Qsq}{\mu^2}.
\end{equation}
O integrando $1/k_T^2$ vem do propagador do quark quase real depois da emissão. O
limite superior é o próprio $\Qsq$: a sonda não resolve emissões mais duras que
ela mesma. O limite inferior $\mu^2$ é onde a perturbação deixa de valer.

O resultado é que a densidade efetiva de quarks vista por uma sonda de
virtualidade $\Qsq$ \emph{depende de $\Qsq$}, e depende logaritmicamente. Não é
uma quebra violenta do escalonamento; é uma deriva lenta. Foi exatamente uma
deriva lenta que se observou.%
\nota{Que a quebra seja logarítmica, e não uma potência, é o que torna QCD testável com os dados de HERA em quatro décadas de $\Qsq$.}
\end{deducao}

\section{As equações DGLAP}

Organizar esses logaritmos em todas as ordens leva às equações de
Dokshitzer--Gribov--Lipatov--Altarelli--Parisi~\cite{Gribov:1972ri,Altarelli:1977zs,Dokshitzer:1977sg}:
\begin{equation}
  \frac{\partial}{\partial \ln \Qsq}
  \begin{pmatrix} q_i(x,\Qsq)\\[2pt] g(x,\Qsq)\end{pmatrix}
  = \frac{\alpha_s(\Qsq)}{2\pi}\int_x^1\frac{\md z}{z}
  \begin{pmatrix} P_{qq}(z) & P_{qg}(z)\\ P_{gq}(z) & P_{gg}(z)\end{pmatrix}
  \begin{pmatrix} q_i(x/z,\Qsq)\\[2pt] g(x/z,\Qsq)\end{pmatrix}.
  \label{eq:dglap}
\end{equation}

Leia a equação em voz alta e ela vira física: \emph{a taxa de variação da
densidade de pártons com $\ln\Qsq$ é dada pela probabilidade de que um párton de
fração maior $x/z$ tenha emitido algo e virado um párton de fração $x$}. As
funções de \emph{splitting} $P_{ab}(z)$ são as probabilidades desses processos
elementares; a mais simples é
\begin{equation}
  P_{qq}(z) = C_F\left[\frac{1+z^2}{(1-z)_+} + \frac{3}{2}\delta(1-z)\right],
  \qquad C_F = \frac{4}{3},
\end{equation}
onde a prescrição-$+$ e o termo com $\delta$ implementam o cancelamento entre
emissões reais e correções virtuais.

\begin{cuidado}[O que DGLAP faz e o que não faz]
DGLAP \emph{não} prevê as PDFs. Ele prevê como elas mudam com $\Qsq$. As PDFs a
uma escala inicial têm de ser extraídas de dados --- é isso que um ``ajuste
global'' faz, e é isso que o HERAPDF2.0 é. Essa distinção importa para o
HADROS3: quando extrapolamos para $x$ pequeníssimo (\cref{cap:cc}), estamos
extrapolando a \emph{condição inicial}, não a evolução. É a parte que a teoria
não controla.
\end{cuidado}

\section{A violação de escala, nos dados}

\begin{figure}[p]
  \centering
  \fig{fig_violacao_escala.pdf}
  \caption{A previsão de DGLAP contra a natureza. Cada série é um $x$ fixo,
  deslocada verticalmente por $2^{i}$ para não se sobrepor. Dados combinados de
  H1 e ZEUS, corrente neutra $e^+p$, $\sqrt{s}=318\gev$, restritos a $y<0{,}6$
  para que $\sigma_r \simeq F_2$~\cite{H1:2015ubc}. Em $x$ pequeno a seção de
  choque \emph{cresce} com $\Qsq$; em torno de $x\approx0{,}13$ ela é plana
  (escalonamento exato!); em $x$ grande ela \emph{decresce}. É exatamente o
  padrão que a \cref{eq:dglap} prevê: aumentar $\Qsq$ resolve emissões, o que
  desloca pártons de $x$ grande para $x$ pequeno.}
  \label{fig:violacao}
\end{figure}

A \cref{fig:violacao} é, na nossa opinião, a figura mais bonita da física de
partículas. Ela mostra três comportamentos qualitativamente distintos num único
gráfico --- crescente, plano, decrescente --- e todos os três saem da mesma
equação, sem ajuste livre no padrão. O ponto de cruzamento, onde o escalonamento
é exato, não é escolhido: é onde o fluxo de pártons para dentro e para fora do
intervalo se cancela.

\section{Liberdade assintótica}

O que torna tudo isso calculável é a liberdade
assintótica~\cite{Gross:1973id,Politzer:1973fx}: a constante de acoplamento forte
decresce com a escala,
\begin{equation}
  \alpha_s(\Qsq) \simeq \frac{1}{b_0 \ln(\Qsq/\Lambda_{\rm QCD}^2)},
  \qquad b_0 = \frac{33-2n_f}{12\pi}.
\end{equation}
É por isso que a hipótese ``pártons livres'' do \cref{cap:partons} funciona: em
$\Qsq$ grande, os pártons realmente são quase livres.%
\nota{$\Lambda_{\rm QCD}\approx 0{,}2$\,GeV é a escala em que $\alpha_s$ explode --- e, não por acaso, $\hbar c/\Lambda_{\rm QCD}\approx 1$\,fm.} E é por isso que a
correção é logarítmica e não catastrófica.

\section{Teste: o modelo funciona?}

\begin{figure}[t]
  \centering
  \fig{fig_f2_modelo_partons.pdf}
  \caption{A \cref{eq:F2em} avaliada com as PDFs HERAPDF2.0 LO, contra a seção de
  choque reduzida medida em $\Qsq\simeq100\gev^2$~\cite{H1:2015ubc}. A curva
  tracejada mostra o mesmo cálculo omitindo charm e bottom: em $x$ pequeno, os
  quarks pesados respondem por quase $30\%$ de $F_2$. Não é um ajuste --- os
  pontos e a curva vêm da mesma publicação, mas a curva é a fórmula do modelo de
  pártons com as PDFs postas dentro.}
  \label{fig:f2partons}
\end{figure}

\begin{cuidado}[Honestidade sobre a \cref{fig:f2partons}]
A comparação acima não é rigorosa, e vale dizer por quê. Primeiro, o que está
plotado como dado é $\sigma_r$, não $F_2$: eles diferem pelos termos de $F_L$ e
$xF_3$, pequenos mas não nulos nessa cinemática. Segundo, $F_2$ foi calculado a
\emph{leading order}, enquanto as PDFs foram ajustadas com uma definição
consistente de esquema. A concordância de poucos por cento que se vê é boa
demais para ser levada a sério em três casas, e ruim demais para ser publicada.
Ela serve para o que serve aqui: mostrar que a estrutura da \cref{eq:F2em} é a
certa.
\end{cuidado}

\section{$x$ pequeno: onde tudo fica difícil}

Em $x \to 0$ a densidade de glúons cresce muito rápido. DGLAP soma logaritmos de
$\Qsq$; em $x$ pequeno também há logaritmos de $1/x$, e somá-los leva à equação
BFKL~\cite{Balitsky:1978ic}. Mas nenhuma equação linear pode valer para sempre:
se a densidade de glúons cresce sem limite, em algum momento eles começam a se
recombinar, e o crescimento \emph{satura}. Essa é a física do Condensado de Vidro
de Cor, descrita por equações não lineares como
Balitsky--Kovchegov~\cite{Balitsky:1995ub,Kovchegov:1999yj}.

Os dois modelos que o HADROS3 tabela vivem aqui:

\begin{description}[style=nextline,leftmargin=0pt,font=\sffamily\bfseries\color{acento}]
\item[GBW (Golec-Biernat--Wüsthoff)] Um modelo de dipolo de cor com saturação
  explícita~\cite{Golec-Biernat:1998zce,Golec-Biernat:1999qor}. O fóton virtual
  flutua num par $q\bar q$ (um ``dipolo'') de tamanho transverso $r$, que espalha
  no próton com uma seção de choque que satura quando $r$ excede
  $1/Q_s(x)$, a escala de saturação.
\item[IIM (Iancu--Itakura--Munier)] Um refinamento que interpola corretamente
  entre o regime de cor transparente e o saturado, ancorado nas soluções
  assintóticas de BK~\cite{Iancu:2003ge}.
\end{description}

% =====================================================================
\chapter{Corrente carregada e neutrinos de energia ultra-alta}
\label{cap:cc}

\section{A massa do $W$ muda tudo}

No caso eletromagnético o propagador do fóton dá $1/(\Qsq)^{2}$ na seção de choque, e
a integral $\int \md \Qsq/(\Qsq)^{2}$ é dominada pelo menor $\Qsq$ possível --- por
isso o espalhamento é predominantemente ``mole''. Para a corrente carregada o
propagador é massivo:
\begin{equation}
  \frac{1}{(\Qsq)^{2}} \;\longrightarrow\;
  \frac{1}{(\Qsq + m_W^2)^2}
  \;=\;\frac{1}{m_W^4}\left(\frac{m_W^2}{m_W^2+\Qsq}\right)^{\!2}.
\end{equation}

\begin{figure}[t]
  \centering
  \fig{fig_propagador.pdf}
  \caption{O fator de propagador. Para $\Qsq \ll m_W^2$ ele é constante --- a
  interação é efetivamente pontual, e recupera-se a teoria de Fermi. Para
  $\Qsq \gg m_W^2$ ele cai como $1/(\Qsq)^{2}$ e sufoca a seção de choque. A
  transição em $\Qsq = m_W^2$ é o que impede $\sigma_{\nu N}$ de crescer
  linearmente para sempre.}
  \label{fig:propagador}
\end{figure}

A consequência é dupla e vale enunciar com clareza:

\begin{enumerate}[leftmargin=*]
\item \textbf{Em baixa energia} ($s \ll m_W^2$), o fator é $\approx1$, a interação
  é pontual, e $\sigma \propto s \propto E_\nu$: a seção de choque cresce
  \emph{linearmente} com a energia. É o regime medido em aceleradores,
  $\sigma/E = 0{,}677\times10^{-38}\cm^2/\gev$ para $\nu$ em alvo
  isoescalar~\cite{Formaggio:2012cpf}.
\item \textbf{Em alta energia} ($s \gg m_W^2$), o propagador impõe
  $\Qsq \lesssim m_W^2$ efetivamente. Junto com $\Qsq = xys$, isso força
  \begin{equation}
    x \;\lesssim\; \frac{m_W^2}{y\,s} \;\sim\; \frac{m_W^2}{2M E_\nu}.
    \label{eq:xtipico}
  \end{equation}
  A seção de choque passa a ser controlada pelas PDFs em $x$ cada vez menor, e o
  crescimento desacelera para $\sigma \propto E^{0{,}36}$.%
\nota{O expoente $0{,}36$ não é fundamental: ele reflete o crescimento das PDFs em $x$ pequeno, e muda se o modelo de saturação mudar.}
\end{enumerate}

\section{Estrutura de sabor}

A corrente carregada não é cega ao sabor. Um $\nu_\ell$ transforma um quark do
tipo \emph{down} num do tipo \emph{up} (ou um antiquark \emph{up} num
\emph{down}), com peso dado pelos elementos da matriz CKM:
\begin{equation}
  F_2^{\rm CC}(x,\Qsq) = 2x\sum_{i,j}\big|V_{q_iq_j}\big|^2\,f_{q_i}(x,\Qsq),
  \qquad q_i\in\{d,s,b\}\ \text{e antiquarks do tipo \emph{up}} .
\end{equation}
Duas consequências práticas: (i) $\nu$ e $\bar\nu$ têm seções de choque
diferentes, com razão $\approx\!2$ em baixa energia, convergindo para 1 em UHE
(porque o mar domina e $q\approx\bar q$); (ii) o canal $s\to c$ é uma fonte
importante de charm, o que tem consequências para a fenomenologia de neutrinos
atmosféricos.

\section{Do keV ao ZeV}

\begin{figure}[t]
  \centering
  \fig{fig_sigma_energia.pdf}
  \caption{A seção de choque $\nu N$ de corrente carregada em treze ordens de
  grandeza de energia. O segmento grosso é o regime linear \emph{medido} em
  aceleradores~\cite{Formaggio:2012cpf}; a linha tracejada é a parametrização de
  Gandhi \emph{et al.}~\cite{Gandhi:1998ri}, na faixa em que os autores a
  declararam válida; as duas curvas cheias são as tabelas GBW e IIM embarcadas no
  HADROS3.}
  \label{fig:sigmaE}
\end{figure}

\begin{cuidado}[Uma discrepância real, ainda não explicada]
Compare, na \cref{fig:sigmaE}, as tabelas do HADROS3 com a referência de
Gandhi \emph{et al.} A tabela GBW fica acima dela por um fator que vai de $\sim7$
em $10^4\gev$ a $\sim90$ em $10^6\gev$; a IIM, por fatores entre $\sim1$ e
$\sim13$. Cálculos modernos com PDFs
ajustadas~\cite{Cooper-Sarkar:2011jtt,Connolly:2011vc,Bertone:2018dse} concordam
entre si dentro de dezenas de por cento nessa faixa, e concordam com a medida do
IceCube~\cite{IceCube:2017roe}. Fatores de dezenas estão fora disso.

A conversão de unidades nas tabelas está correta (conferimos: a coluna em
$\gev^{-2}$ vezes $3{,}894\times10^{-28}\cm^2\gev^{2}$ reproduz a coluna em
$\cm^2$). Logo a diferença está no conteúdo, não na unidade. As hipóteses
plausíveis --- normalização por núcleo em vez de por núcleon, soma de CC${}+{}$NC,
ou um problema na geração das tabelas --- ainda não foram testadas. Como
$\tau \propto \sigma$, um fator 90 aqui é um fator 90 na profundidade óptica.
Esta é uma pendência aberta do HADROS3, não um resultado.
\end{cuidado}

\section{Por que isto é uma extrapolação}

A \cref{eq:xtipico} e a \cref{fig:plano} contam juntas a história incômoda da
área. Um neutrino de $10^{9}\gev$ sonda $x \sim 3\times10^{-6}$; o menor $x$ jamais
medido em DIS foi $\sim5\times10^{-6}$, e apenas em $\Qsq$ muito baixo, longe de
$m_W^2$. Ou seja: \textbf{a seção de choque de neutrinos UHE não é medida, é
extrapolada}, e o que se extrapola é justamente a região em que a teoria linear
deixa de valer.

É por isso que existem modelos de saturação, e é por isso que faz diferença qual
deles se usa. As previsões modernas diferem entre si por dezenas de por cento até
$10^{12}\gev$~\cite{Cooper-Sarkar:2011jtt,Connolly:2011vc,Bertone:2018dse}; abaixo
disso, o IceCube já mediu a seção de choque em torno de
$10^{4}$--$10^{6}\gev$ usando a absorção pela própria
Terra~\cite{IceCube:2017roe} --- exatamente o cálculo de profundidade óptica do
próximo parágrafo, usado ao contrário. O Electron-Ion Collider deve reduzir parte
dessa incerteza medindo $x$ pequeno com precisão~\cite{Accardi:2012qut,AbdulKhalek:2021gbh}.

\section{Fechando o círculo: a profundidade óptica}

Tudo o que fizemos até aqui produz um único número por energia,
$\sigma_{\nu N}(E)$. É esse número que entra no HADROS3, através da profundidade
óptica:
\begin{equation}
  \tau(E) = \int n_b(\ell)\;\sigma_{\nu N}(E)\;\md\ell
          = \frac{\sigma_{\nu N}(E)}{m_b}\int \rho\,\md\ell,
  \qquad P_{\rm int} = 1 - e^{-\tau}.
  \label{eq:tau}
\end{equation}

\begin{figure}[t]
  \centering
  \fig{fig_profundidade_optica.pdf}
  \caption{A \cref{eq:tau} para três colunas de matéria, usando a tabela GBW.
  A Terra fica opaca em torno de algumas dezenas de TeV --- e é exatamente essa
  opacidade que o IceCube usou para medir $\sigma_{\nu N}$~\cite{IceCube:2017roe}.
  O toro do HADROS3, com a configuração padrão, é opaco por nove ordens de
  grandeza no plano equatorial.}
  \label{fig:tau}
\end{figure}

A separação da \cref{eq:tau} é limpa e vale memorizar: \textbf{toda a física de
partículas está em $\sigma$; toda a astrofísica e geometria estão na coluna
$X=\int\rho\,\md\ell$}. Elas só se encontram no produto. O notebook
\texttt{01\_profundidade\_optica\_dis.ipynb} do sandbox constrói essa integral
com todo o cuidado --- discretização, convergência, e o que a Relatividade Geral
acrescenta quando o meio está girando perto de um buraco negro.

% =====================================================================
\chapter{Para onde ir depois}

\section{Um roteiro de leitura}

\begin{description}[style=nextline,leftmargin=0pt,font=\sffamily\bfseries\color{acento}]
\item[Para consolidar o básico] Halzen \& Martin~\cite{Halzen:1984mc}, capítulos
  8--9, é o tratamento clássico e ainda o mais legível. Thomson~\cite{Thomson:2013zua}
  é mais moderno e mais gentil. Close~\cite{Close:1979bt} é antigo mas
  insuperável na intuição do modelo de pártons.
\item[Para a QCD de verdade] Ellis, Stirling \& Webber~\cite{Ellis:1996mzs} é o
  padrão da área. Para a estrutura formal (fatoração, esquemas, o que
  \emph{significa} uma PDF), Collins~\cite{Collins:2011zzd}.
\item[Para DIS especificamente] Devenish \& Cooper-Sarkar~\cite{Devenish:2004pb}
  é o livro dedicado ao assunto, escrito por gente do HERA;
  Roberts~\cite{Roberts:1990ww} cobre a teoria com mais detalhe técnico.
\item[Para teoria de campos por baixo] Peskin \& Schroeder~\cite{Peskin:1995ev},
  capítulo 17, faz o DIS dentro do formalismo completo; Schwartz~\cite{Schwartz:2014sze}
  é a alternativa mais moderna.
\item[Para os dados] O \emph{Review of Particle Physics}~\cite{ParticleDataGroup:2024cfk}
  tem uma revisão de funções de estrutura curta e confiável, atualizada a cada
  dois anos. E leia o artigo da combinação de HERA~\cite{H1:2015ubc}: são os
  dados que usamos aqui.
\item[Para neutrinos UHE] Gandhi \emph{et al.}~\cite{Gandhi:1998ri} é o ponto de
  partida histórico; Cooper-Sarkar \emph{et al.}~\cite{Cooper-Sarkar:2011jtt},
  Connolly \emph{et al.}~\cite{Connolly:2011vc} e Bertone
  \emph{et al.}~\cite{Bertone:2018dse} são as referências modernas.
  Formaggio \& Zeller~\cite{Formaggio:2012cpf} compila tudo o que foi medido.
\end{description}

\section{Exercícios}

\begin{enumerate}[leftmargin=*,itemsep=5pt]
\item Mostre que $q\cdot k = \Qsq/2$ e $q\cdot k' = -\Qsq/2$ no limite de massas
  leptônicas nulas, completando a verificação de $q_\mu L^{\mu\nu}=0$ da
  \cref{eq:Lem}.

\item Refaça a \cref{eq:W2x} \emph{sem} desprezar a massa do núcleon nem a do
  lépton. Em que região do plano $(x,\Qsq)$ as correções de massa passam de
  $1\%$? Compare com a cobertura dos dados na \cref{fig:plano}.

\item Complete a Dedução 3.3: escreva explicitamente as condições
  $q_\mu W^{\mu\nu}=0$ para a \cref{eq:Wgeral} e obtenha $W_4$, $W_5$, $W_6$ em
  termos de $W_1$ e $W_2$.

\item Deduza a distribuição em $y$ do processo $\bar\nu q$ e do $\bar\nu\bar q$,
  e explique por que $\sigma_{\bar\nu N}/\sigma_{\nu N} \to 1$ quando o mar
  domina.

\item A partir das PDFs em \texttt{sandbox/data/hera/herapdf20\_lo\_xf.dat},
  calcule a fração do momento do próton carregada pelos quarks e a carregada
  pelo glúon, em $\Qsq = 10\gev^2$ e $10^4\gev^2$. Confirme que a soma dá 1 e
  que a partilha muda com a escala.

\item Estime, com a \cref{eq:xtipico}, o $x$ típico sondado por um neutrino de
  $6{,}3\,$PeV --- a energia da ressonância de Glashow. Você acha que uma
  extrapolação de PDF até lá é confiável? Justifique com a \cref{fig:plano}.

\item \textbf{(Aberto.)} Investigue a discrepância documentada na
  \cref{fig:sigmaE}: as tabelas do HADROS3 são por núcleon ou por núcleo? Somam
  CC${}+{}$NC? Comece por \texttt{cpp/apps/hadros3\_dis\_sampler.cpp} e pela
  geração das tabelas em \texttt{data/sigma/}. A resposta não está nesta
  apostila.
\end{enumerate}

% =====================================================================
\cleardoublepage
\phantomsection
\addcontentsline{toc}{chapter}{Referências}
\printbibliography[title={Referências}]

\end{document}
